\begin{lemma}
Let \(C=100\), let \(\beta=1/2\), and put
\[
M(n)=\noderef{TwoBiteNaturalReducedVertexCount}(n)
=\left\lfloor n/\log^2 n\right\rfloor,\qquad
p(n)=\noderef{TwoBiteEdgeProbability}(\beta,n)
=\frac12\sqrt{\frac{\log n}{n}} .
\]
Under the two-bite law with base size \(M(n)\) and edge probability \(p(n)\),
with high probability the following conditional collision statement holds.
For every sample \(\omega\), if every lifted neighborhood satisfies
\[
|N_x|\le 2p(n)n\qquad (x\in V_R\cup V_B),
\tag{1}
\]
then for every two distinct red vertices \(r\ne s\),
\[
|\pi_B(N_r)\cap\pi_B(N_s)|
\le |N_r\cap N_s|+C(\log n)^3,
\tag{2}
\]
and for every two distinct blue vertices \(b\ne c\),
\[
|\pi_R(N_b)\cap\pi_R(N_c)|
\le |N_b\cap N_c|+C(\log n)^3.
\tag{3}
\]
\end{lemma}
\begin{proof}
Write \(L=\log n\), \(M=M(n)\), and \(p=p(n)\).  We prove that the
probability that the implication in the statement fails is \(o(1)\).  It
suffices to prove the red estimate (2), because the proof of (3) is obtained
by interchanging the two colors and replacing blue projections by red
projections.

Fix two distinct red vertices \(r\ne s\).  Let
\[
A=N_{G_R}(r),\qquad B=N_{G_R}(s),
\]
and put
\[
U=N_r,\qquad V=N_s,\qquad U_0=U\setminus V,\qquad V_0=V\setminus U.
\]
The only random coordinates relevant to the red estimate are the red base
graph and the injection
\(\pi=\noderef{TwoBiteEmbedding}(\omega)\); the blue base graph may be summed
out because its \(\noderef{GnpGraphWeight}\) total mass is \(1\).
Condition on the red base graph and on the complete red-coordinate map
\[
\rho(v)=\pi_R(v),\qquad v\in V(G).
\]
This conditioning fixes \(A,B,U,V,U_0,V_0\).  If the size hypothesis (1) fails,
the implication in the statement is true, so in bounding failure we may assume
\[
a:=|U_0|\le |U|\le 2pn,\qquad
b:=|V_0|\le |V|\le 2pn.
\tag{4}
\]

We now spell out the remaining conditional law from
\(\noderef{UniformInjectionWeight}\).  For each red row \(i\in V_R\), let
\[
R_i=\{v:\rho(v)=i\},\qquad q_i=|R_i|.
\]
If some \(q_i>M\), then no injection with red-coordinate map \(\rho\) exists
and this conditional fiber has probability \(0\).  Otherwise the number of
injection completions with this red-coordinate map is
\[
\prod_{i\in V_R} (M)_{q_i},
\]
where \((M)_q=M(M-1)\cdots(M-q+1)\).  Indeed, in row \(i\) one must assign
distinct blue coordinates to the \(q_i\) vertices of \(R_i\), and choices in
different red rows are independent because different red coordinates already
make the product cells distinct.  Since
\(\noderef{UniformInjectionWeight}\) is constant over all injections and
\(\noderef{TwoBiteSampleWeight}\) is the product of the two graph weights with
this uniform-injection factor, conditioning on the nonzero fiber makes the
blue-coordinate assignments in the rows \(R_i\) independent, each uniformly
distributed over injections \(R_i\hookrightarrow V_B\).

The row sets involved in \(U_0\) and \(V_0\) are disjoint: \(U_0\) uses red
rows in \(A\setminus B\), while \(V_0\) uses red rows in \(B\setminus A\).
Therefore, after conditioning further on all blue coordinates of vertices in
\(U_0\), the assignments on the rows meeting \(V_0\) are still independent
uniform row injections.  Let
\[
S=\pi_B(U_0)\subseteq V_B.
\]
Then \(|S|\le a\).  Define
\[
W=\bigl|\{v\in V_0:\pi_B(v)\in S\}\bigr|.
\]
Every blue value in \(\pi_B(U_0)\cap\pi_B(V_0)\) has at least one witnessing
vertex in \(V_0\) whose blue coordinate lies in \(S\), hence
\[
|\pi_B(U_0)\cap\pi_B(V_0)|\le W. \tag{5}
\]

We bound the tail of \(W\) by a direct counting argument.  Fix an integer
\(t\ge1\).  If \(W\ge t\), then some \(t\)-element set
\(T\subseteq V_0\) has all of its blue coordinates in \(S\).  For a fixed
such \(T\), write \(h_i=|T\cap R_i|\) for each red row used by \(V_0\).  Under
the row-injection law, the probability that all \(h_i\) specified vertices in
row \(i\) land in \(S\) is
\[
\frac{(|S|)_{h_i}}{(M)_{h_i}}
\le \left(\frac{|S|}{M}\right)^{h_i}
\le \left(\frac{a}{M}\right)^{h_i}.
\]
The first inequality follows term by term:
\[
\frac{|S|-j}{M-j}\le \frac{|S|}{M}\qquad (0\le j<h_i),
\]
because \(|S|\le M\).  Multiplying over the rows, which are conditionally
independent, gives
\[
\Pr(T\hbox{ lands in }S\mid \rho,G_R,S)
\le \left(\frac{a}{M}\right)^t .
\]
There are at most \(\binom bt\) choices for \(T\).  Hence
\[
\Pr(W\ge t\mid \rho,G_R,S)
\le \binom bt\left(\frac{a}{M}\right)^t
\le \left(\frac{eab}{Mt}\right)^t . \tag{6}
\]
If \(t>b\), then \(W\ge t\) is impossible, so the same displayed upper bound
is harmless; otherwise the last inequality is the standard
\(\binom bt\le(eb/t)^t\).
This estimate is uniform in the conditioning data, so it remains valid after
averaging over \(S,\rho\), and \(G_R\).

On the size event (4),
\[
\frac{ab}{M}
\le \frac{(2pn)^2}{M}
= \frac{4p^2n^2}{M}.
\]
Since \(p^2n=L/4\) and, for all sufficiently large \(n\),
\[
M=\left\lfloor\frac n{L^2}\right\rfloor
\ge \frac12\,\frac n{L^2},
\]
we get
\[
\frac{ab}{M}\le 2L^3. \tag{7}
\]
Take \(t=\lceil C L^3\rceil\).  Combining (6) and (7), and increasing \(n\)
if necessary so that \(L\ge1\), gives
\[
\Pr(W\ge C L^3\mid r,s,\hbox{ size event})
\le
\left(\frac{2eL^3}{C L^3}\right)^{C L^3}
\le \exp(-cL^3)
\tag{8}
\]
for an absolute constant \(c>0\), because \(C=100>2e\).

It remains only to relate \(W\) to the projected intersection in (2).  If
\(y\in\pi_B(U)\cap\pi_B(V)\), choose \(u\in U\) and \(v\in V\) with
\(\pi_B(u)=y=\pi_B(v)\).  If \(u\in V\) or \(v\in U\), then
\(y\in\pi_B(U\cap V)\).  Otherwise \(u\in U_0\) and \(v\in V_0\), so
\(y\in\pi_B(U_0)\cap\pi_B(V_0)\).  Thus
\[
\pi_B(U)\cap\pi_B(V)
\subseteq
\pi_B(U\cap V)\cup(\pi_B(U_0)\cap\pi_B(V_0)).
\]
Taking cardinalities and using \(|\pi_B(U\cap V)|\le |U\cap V|\) and (5),
we obtain
\[
|\pi_B(U)\cap\pi_B(V)|
\le |U\cap V|+W. \tag{9}
\]
Equations (8) and (9) show that, for this fixed red pair, the probability of
violating (2) while the size hypothesis holds is at most \(\exp(-cL^3)\).

There are fewer than \(M^2\le n^2=\exp(2L)\) ordered red pairs.  The union
bound gives a total red-pair failure probability at most
\[
M^2\exp(-cL^3)=o(1).
\]
Repeating the same argument with colors swapped gives the same bound for all
blue pairs.  Therefore the probability that either (2) or (3) fails for some
same-color pair while (1) holds is \(o(1)\).  In the tablet support DAG, the
red and blue one-color high-probability statements are packaged as
\(\noderef{OppositeProjectionRedFailureBound}\) and
\(\noderef{OppositeProjectionSymmetricBlue}\).  Their intersection is obtained
using \(\noderef{WithHighProbabilityThreeIntersection}\), with the red event
repeated as the third event and with eventual nonnegativity of sample weights
coming from \(\noderef{OppositeProjectionEdgeProbBounds}\) and
\(\noderef{TwoBiteSampleWeightNonnegative}\).  The resulting triple
intersection implies the desired two-event conjunction.  This is exactly the
\(\noderef{WithHighProbability}\) assertion for the conditional collision
statement under the two-bite probability
\(\noderef{TwoBiteEventProbability}\), with the rounded base size
\(M(n)=\noderef{TwoBiteNaturalReducedVertexCount}(n)\) and
\(p(n)=\noderef{TwoBiteEdgeProbability}(1/2,n)\).
\end{proof}
